\documentclass[report.tex]{subfiles}

\begin{document}

\section{Coinduction Example: Stream Bisimilarity}
\label{app:stream-bisimilarity}

To illustrate how coinductive reasoning operates in practice,
consider the \emph{bisimilarity} relation $\approx$ on streams of natural
numbers,
$\Stream{\mathbb{N}}$,
defined coinductively by the following inference rule:
\begin{equation*}
  \begin{prooftree}
    \hypo{n_1 = n_2}
    \hypo{s_1 \approx s_2}
    \infer[double]2[Bisim-Step]{
      \mathsf{scons}(n_1, s_1) \approx \mathsf{scons}(n_2, s_2)}
  \end{prooftree}
\end{equation*}

Let us prove that this relation is reflexive:
\begin{equation*}
  \forall s \in \Stream{\mathbb{N}}, s \approx s.
\end{equation*}
The proof proceeds by coinduction on the \emph{bisimilarity} relation:
\begin{enumerate}
  \item Applying the coinduction principle provides the coinductive
        hypothesis:
        \begin{equation*}
          \textsf{CoIH} : \forall s' \in \Stream{\mathbb{N}}, s' \approx s',
        \end{equation*}
        with the restriction that $\textsf{CoIH}$ may only be applied under a
        constructor of the coinductive relation.
  \item We decompose the stream $s$ into a head $n \in \mathbb{N}$ and a
        tail $s' \in \Stream{\mathbb{N}}$,
        so that $s = \mathsf{scons}(n, s')$.
        The goal becomes:
        \begin{equation*}
          \mathsf{scons}(n, s') \approx \mathsf{scons}(n, s').
        \end{equation*}
  \item We apply the coinductive rule \textsf{Bisim-Step}.
        The first premise $n = n$ holds trivially,
        leaving the second premise $s' \approx s'$ to be established.
  \item Because the constructor Bisim-Step has been applied,
        we can now legitimately invoke $\textsf{CoIH}$.
        This immediately concludes the proof.
\end{enumerate}


The complete derivation tree for the loop iteration case ($b = \mathtt{true}$)
is displayed in \cref{fig:free-sim-example-derivation}.

\begin{figure}[htb]
  \centering
  \begin{equation*}
  \begin{prooftree}
    \hypo{0 < 2}
    \hypo{0 < 4}
    \infer2[\textsf{CIH}]{\fsim{P_t}{2}{4}{P_s}{=}}
    \infer1[\textsc{FSourceSteps}\;($\times 3$)]{%
      \fsim{P_t}{2}{1}{\mathtt{skip};\mathtt{skip};\mathtt{skip};P_s}{=}}
    \infer1[\textsc{FSourceSteps}]{\fsim{P_t}{2}{0}{P_s}{=}}
    \infer1[\textsc{FTargetSteps}]{\fsim{\mathtt{skip};P_t}{1}{0}{P_s}{=}}
    \infer1[\textsc{FTargetSteps}]{\fsim{P_t}{0}{0}{P_s}{=}}
  \end{prooftree}
\end{equation*}
\caption{Derivation tree for the loop iteration case ($b = \mathtt{true}$).}
  \label{fig:free-sim-example-derivation}
\end{figure}


% Appendix ideam details the definition of meminj and its lemmas
\end{document}

% Local Variables:
% TeX-engine: luatex
% End:

% LocalWords:  Coinductive coinductive asymmetricity arities
